\section*{Chapter \ref{ch:b1:carbohydrates} --- Carbohydrates}

\begin{solution}{exo:b1:carbohydrates:1}
Glucose: hexose, \omterm{def:b1:carbohydrates:monosaccharide}{aldose}. Fructose: hexose, \omterm{def:b1:carbohydrates:monosaccharide}{ketose}. Ribose: pentose,
\omterm{def:b1:carbohydrates:monosaccharide}{aldose}. Glyceraldehyde: triose, \omterm{def:b1:carbohydrates:monosaccharide}{aldose}.
\end{solution}

\begin{solution}{exo:b1:carbohydrates:2}
The carbonyl carbon after ring closure, bearing the ring oxygen and a
hydroxyl, which can open back to the aldehyde or ketone. In sucrose
both anomeric carbons are locked in the \omterm{def:b1:carbohydrates:glycosidic}{glycosidic bond}; neither can
open, so no aldehyde forms and the copper is not reduced.
\end{solution}

\begin{solution}{exo:b1:carbohydrates:3}
\omterm{def:b1:carbohydrates:polysaccharide}{Starch}: glucose, $\alpha(1{\to}4)$ with $\alpha(1{\to}6)$ branches
(amylopectin). \omterm{def:b1:carbohydrates:polysaccharide}{Glycogen}: the same, more branched. \omterm{def:b1:carbohydrates:polysaccharide}{Cellulose}: glucose,
$\beta(1{\to}4)$. \omterm{def:b1:carbohydrates:polysaccharide}{Chitin}: N-acetylglucosamine, $\beta(1{\to}4)$.
\end{solution}

\begin{solution}{exo:b1:carbohydrates:4}
C1 down, C2 down, C3 up, C4 down. $\alpha$-D-galactose is the same with
the C4 hydroxyl up.
\end{solution}

\begin{solution}{exo:b1:carbohydrates:5}
In the $\alpha$ bond the anomeric hydroxyl points down, and each
residue can be added at the same angle as the last: the chain turns
steadily and closes into a helix. In the $\beta$ bond it points up, and
a residue can only be added flipped over: the chain runs straight.
Iodine molecules fit inside the amylose helix and their electronic
state changes (blue); a flat ribbon offers no cavity.
\end{solution}

\begin{solution}{exo:b1:carbohydrates:6}
Iodine: only \omterm{def:b1:carbohydrates:polysaccharide}{starch} turns blue-black. Benedict on the other two: only
glucose gives a red precipitate; sucrose stays blue. Confirm sucrose by
heating with dilute acid, neutralising, and repeating Benedict: now
positive.
\end{solution}

\begin{solution}{exo:b1:carbohydrates:7}
Chains of 12 residues: about \num{4000} chains, half of them outermost:
some \num{2000} ends. Amylose: one end. To release \num{5000} glucoses
at one per second per end: \omterm{def:b1:carbohydrates:polysaccharide}{glycogen} \qty{2.5}{s}; amylose \qty{5000}{s},
nearly an hour and a half.
\end{solution}

\begin{solution}{exo:b1:carbohydrates:8}
$\sigma = 0.8\times 25/(2\times 0.5) = \qty{20}{MPa}$: a fiftieth of the
fibril strength; the wall has a large safety margin, needed because the
fibrils are only a quarter of the wall.
\end{solution}

\begin{solution}{exo:b1:carbohydrates:9}
Lactase breaks the $\beta(1{\to}4)$ bond between galactose and glucose;
without it lactose is not absorbed (only \omterm{def:b1:carbohydrates:monosaccharide}{monosaccharides} cross the
\omterm{def:b1:body-plans-tissues:epithelium}{epithelium}). In the colon bacteria ferment it to acids and gas
(cramps, bloating), and the undigested sugar and acids draw water
osmotically into the lumen: diarrhoea.
\end{solution}

\begin{solution}{exo:b1:carbohydrates:10}
Bacteria and protists in the rumen secrete cellulases and ferment the
glucose to short \omterm{def:b1:lipids:fattyacid}{fatty acids}, which the cow absorbs and oxidises; the
microbes themselves are later digested. The fermentation must precede
the small intestine so that its products reach the absorbing surface
and the microbial protein reaches the stomach and intestine. \omterm{def:b1:carbohydrates:polysaccharide}{Cellulose}:
\qty{4}{kg}; fermented \qty{2.4}{kg}; \qty{29}{MJ} a day, most of the
cow's energy.
\end{solution}

\begin{solution}{exo:b1:carbohydrates:11}
O \omterm{def:b1:cell-unit-of-life:cell}{cells} carry neither antigen and provoke no antibody in anyone; O
plasma, however, holds antibodies against both A and B, so an O person
attacks any A or B \omterm{def:b1:cell-unit-of-life:cell}{cell} received. An A person's immune system has never
seen B and makes anti-B antibodies (against gut bacteria bearing
similar sugars), but tolerates A: A plasma contains anti-B only.
\end{solution}

\begin{solution}{exo:b1:carbohydrates:12}
Right: glucose is the fuel of nearly every \omterm{def:b1:cell-unit-of-life:cell}{cell} and \omterm{def:b1:carbohydrates:polysaccharide}{cellulose} the most
abundant organic molecule on Earth, and both are polymers of, or the
same as, D-glucose. One bond changes everything: $\alpha$ gives a coil
that enzymes open and \omterm{def:b1:cell-unit-of-life:cell}{cells} burn; $\beta$ gives a fibre that resists
enzymes, water and tension, and that most animals cannot use at all.
Biology is not the list of molecules but the geometry of their
joining --- the same chemistry can be food or scaffold according to a
single stereochemical choice, and \omterm{def:b1:organism-environment:organism}{organisms} are separated by which
enzymes they possess to undo it.
\end{solution}

\begin{solution}{pb:b1:carbohydrates:1}
\textbf{1.} $2^{t-1}$: 1, 2, 4, and $2^{11} = 2048$.
\textbf{2.} $2^{12} - 1 = 4095$.
\textbf{3.} $4095\times 13 = \num{53235}$, about \num{53000}.
\textbf{4.} $\num{53235}\times 162 = \qty{8.6e6}{Da}$;
$\qty{1.43e-17}{g}$.
\textbf{5.} $2048/4095$: half. \num{2048} non-reducing ends.
\textbf{6.} One.
\textbf{7.} $12\times 1.9 = \qty{23}{nm}$ radius, \qty{46}{nm}
diameter: twice a ribosome.
\textbf{8.} $100/\num{1.43e-17} = \num{7.0e18}$ particles.
\textbf{9.} $100/162 = \qty{0.62}{mol}$ of residues, i.e.\ \qty{111}{g}
of free glucose (the water of \omterm{prop:b1:water-small-molecules:families}{hydrolysis} adds the difference).
\textbf{10.} \qty{400}{g}, over a quarter of the liver's mass.
\textbf{11.} \qty{10}{g/h} $= \qty{2.78}{mg/s} = \qty{1.54e-5}{mol/s} =
\num{9.3e18}$ molecules per second.
\textbf{12.} About \qty{11}{h} (111 g at 10 g/h) --- a night. Chiefly
the brain, which takes \qty{5}{g/h}.
\textbf{13.} $2048\times\num{7e18} = \num{1.4e22}$ ends.
\textbf{14.} $\num{1.4e22}\times 20 = \num{2.9e23}$ per second $=
\qty{0.48}{mol/s} = \qty{86}{g/s}$: thirty thousand times the actual
rate. The ends are never limiting; the enzyme is.
\textbf{15.} $\num{9.3e18}/20 = \num{4.6e17}$ molecules --- about
\qty{75}{mg} of enzyme, a small fraction of the liver's protein.
\textbf{16.} Chains of \num{53000} residues: still $\num{7e18}$ ends
(one each), \num{2048} times fewer; maximum rate \qty{42}{mg/s}, still
fifteen times the actual rate. (With longer amylose chains the margin
shrinks; the real limit of an unbranched store is its insolubility and
crystallisation, which take the ends out of reach.)
\textbf{17.} When blood glucose falls, hormones switch phosphorylase
on within seconds; the enzyme finds thousands of ends per particle
already exposed at the surface, so the release rate can rise a
hundredfold at once, without waiting for chains to unwind or
particles to dissolve.
\textbf{18.} $0.55/1.0 = \qty{0.55}{osmol/L}$ (\qty{100}{g}
$= \qty{0.55}{mol}$ of free glucose).
\textbf{19.} Nearly double the \omterm{def:b1:eukaryotic-cell:organelle}{cytosol}'s osmolarity: water would rush
in from the blood, the \omterm{def:b1:cell-unit-of-life:cell}{cells} would swell to almost three times their
volume and burst.
\textbf{20.} $\num{7e18}/\num{6e23} = \qty{1.2e-5}{mol}$ in \qty{1}{L}:
\qty{12}{\micro osmol/L}, forty thousand times less.
\textbf{21.} $\Psi_s = -2.58\times 0.55 = \qty{-1.4}{MPa}$: a wall
would need to hold \qty{1.4}{MPa}, fourteen atmospheres.
\textbf{22.} Muscle \omterm{def:b1:carbohydrates:polysaccharide}{glycogen} is the muscle's own fuel for contraction,
mobilised as glucose-6-phosphate straight into glycolysis; the muscle
lacks the enzyme that frees glucose from its phosphate, so the store
cannot leave the fibre.
\textbf{23.} A potato mobilises its \omterm{def:b1:carbohydrates:polysaccharide}{starch} over weeks of sprouting; a
liver must respond in minutes. Large dense grains with few surface
ends suit slow steady release; small highly branched particles with
thousands of ends suit a fast one.
\textbf{24.} $0.62/(4\times 3600) = \qty{4.3e-5}{mol/s} = \num{2.6e19}$
residues per second; \qty{1.24}{mol} of \omterm{def:b1:water-small-molecules:atp}{ATP} in all, $\num{5.2e19}$ per
second.
\textbf{25.} Release: $\num{9.3e18}$ glucose molecules per second
(\qty{10}{g/h}); made possible by \num{1.4e22} non-reducing ends,
\num{2048} per particle; the polymer spares the \omterm{def:b1:cell-unit-of-life:cell}{cells}
\qty{0.55}{osmol/L}, nearly twice their own osmolarity.
\end{solution}
